\section*{Capítulo \ref{ch:g11:titration} --- Las valoraciones}

\begin{solution}{exo:g11:titration:1}
El \omterm{def:g11:titration:titration}{valorante} es la disolución de concentración conocida con precisión,
en la bureta; la \omterm{def:g11:titration:titration}{disolución valorada} contiene la especie que se quiere
medir, en el matraz. En la \omterm{def:g11:titration:equivalence}{equivalencia}, el \omterm{def:g11:titration:titration}{valorante} añadido y la
especie valorada están en las proporciones de la ecuación: los dos se
han consumido.
\end{solution}

\begin{solution}{exo:g11:titration:2}
$n = 0.0200 \times 0.0125 = \qty{2.50e-4}{mol}$.
\end{solution}

\begin{solution}{exo:g11:titration:3}
$V_E = \qty{14.3}{mL}$; la línea del hierro(II) empieza en
\qty{14.3e-4}{mol}, es decir, \qty{1.43e-3}{mol}.
\end{solution}

\begin{solution}{exo:g11:titration:4}
El propio \omterm{def:g11:titration:titration}{valorante} tiene color (violeta) y su producto es casi
incoloro: el color aparece en el matraz en cuanto el permanganato está
en exceso.
\end{solution}

\begin{solution}{exo:g11:titration:5}
5: la ecuación consume 5 \omterm{def:g9:ions:ion}{iones} hierro(II) por cada \omterm{def:g9:ions:ion}{ion} permanganato, y
en la \omterm{def:g11:titration:equivalence}{equivalencia} los \omterm{def:g7:chemical-reactions:reactant}{reactivos} se han juntado en esas proporciones.
\end{solution}

\begin{solution}{exo:g11:titration:6}
$n(\ce{MnO4^-}) = 0.0200 \times 0.00840 = \qty{1.68e-4}{mol}$;
$n(\ce{Fe^{2+}}) = 5 \times \num{1.68e-4} = \qty{8.40e-4}{mol}$;
$c = \num{8.40e-4} / 0.0100 = \qty{0.0840}{mol/L}$;
$C_m = 0.0840 \times 55.8 = \qty{4.69}{g/L}$.
\end{solution}

\begin{solution}{exo:g11:titration:7}
Razón 1:1: $n = 0.0100 \times 0.0114 = \qty{1.14e-4}{mol}$ en
\qty{10.0}{mL}, así que \qty{1.14e-3}{mol} en los \qty{100.0}{mL}, es
decir, $\num{1.14e-3} \times 176.0 = \qty{0.201}{g}$: unos
\qty{200}{mg}.
\end{solution}

\begin{solution}{exo:g11:titration:8}
Cada gota no reaccionaría enseguida: el color del \omterm{def:g11:titration:titration}{valorante} persistiría
un rato antes de la \omterm{def:g11:titration:equivalence}{equivalencia}, y el final se vería demasiado pronto
(o habría que esperar minutos después de cada gota).
\end{solution}

\begin{solution}{exo:g11:titration:9}
Demasiado grande: se ha pasado de la \omterm{def:g11:titration:equivalence}{equivalencia} y ha añadido
permanganato en exceso. Cerca de la \omterm{def:g11:titration:equivalence}{equivalencia} debería haber añadido
el \omterm{def:g11:titration:titration}{valorante} gota a gota, agitando, y haberse detenido en el primer rosa
tenue que persiste.
\end{solution}

\begin{solution}{exo:g11:titration:10}
$n(\ce{Fe^{2+}}) = 5 \times 0.0200 \times 0.0130 = \qty{1.30e-3}{mol}$
en \qty{10.0}{mL}: $c = \qty{0.130}{mol/L}$ para la disolución diluida,
y $10 \times 0.130 = \qty{1.30}{mol/L}$ para la concentrada.
\end{solution}

\begin{solution}{exo:g11:titration:11}
$0.0630 \times 55.8 = \qty{3.52}{g}$.
\end{solution}

\begin{solution}{exo:g11:titration:12}
$0.1 / 7.2 \approx \qty{1.4}{\percent}$, frente al \qty{0.7}{\percent}
para \qty{14.3}{mL}. Valorar un volumen mayor (o una muestra más
concentrada) de la disolución, o usar un \omterm{def:g11:titration:titration}{valorante} más diluido.
\end{solution}

\begin{solution}{exo:g11:titration:13}
$n(\ce{MnO4^-}) = 0.0200 \times 0.0176 = \qty{3.52e-4}{mol}$;
$n(\ce{H2O2}) = \frac{5}{2} \times \num{3.52e-4} = \qty{8.80e-4}{mol}$ en
\qty{10.0}{mL}: \qty{0.0880}{mol/L} en la disolución diluida, y
\qty{0.880}{mol/L} en el desinfectante, es decir,
$0.880 \times 34.0 = \qty{29.9}{g/L}$ (una disolución al 3 \%, más o
menos).
\end{solution}

\begin{solution}{exo:g11:titration:14}
$n(\ce{MnO4^-}) = \num{1.4e-3} / 5 = \qty{2.8e-4}{mol}$, que hay que
verter en unos \qty{0.015}{L}: $c \approx \qty{0.019}{mol/L}$, así que
una disolución de \qty{0.0200}{mol/L}. Diez veces más concentrada, $V_E$
sería de unos \qty{1.4}{mL}: el error de la bureta llegaría a un
\qty{7}{\percent}.
\end{solution}

\begin{solution}{exo:g11:titration:15}
$8 \times 0.0200 \times 0.0143 = \qty{2.29e-3}{mol}$ de \omterm{def:g9:ions:ion}{iones} hidrógeno.
\qty{10}{mL} a \qty{1.0}{mol/L} contienen \qty{1.0e-2}{mol}: bastan, y
sobran más de cuatro veces. Sin ácido, la reacción del capítulo no
podría producirse tal como está escrita (los \omterm{def:g9:ions:ion}{iones} hidrógeno son un
\omterm{def:g7:chemical-reactions:reactant}{reactivo}), y el permanganato reaccionaría de otra manera.
\end{solution}

\begin{solution}{pb:g11:titration:1}
\textbf{1.} \ce{MnO4^-}/\ce{Mn^{2+}} y \ce{Fe^{3+}}/\ce{Fe^{2+}}. El
\omterm{def:g11:redox:oxidant}{oxidante} es el \omterm{def:g9:ions:ion}{ion} permanganato, y el \omterm{def:g11:redox:oxidant}{reductor}, el \omterm{def:g9:ions:ion}{ion} hierro(II).

\textbf{2.} \ce{MnO4^- + 8H+ + 5e- -> Mn^{2+} + 4H2O};
\ce{Fe^{2+} -> Fe^{3+} + e-}.

\textbf{3.} La segunda se multiplica por 5:
\ce{MnO4^- + 8H+ + 5Fe^{2+} -> Mn^{2+} + 5Fe^{3+} + 4H2O}. Carga a la
izquierda: $-1 + 8 + 10 = +17$; a la derecha: $+2 + 15 = +17$.

\textbf{4.} Los \omterm{def:g9:ions:ion}{iones} hidrógeno son un \omterm{def:g7:chemical-reactions:reactant}{reactivo}: la reacción necesita
una \omterm{def:g9:acids-bases-ph:acidic}{disolución ácida}.

\textbf{5.} Es total, rápida y la única reacción del permanganato en el
matraz; el permanganato, violeta, se consume antes de la \omterm{def:g11:titration:equivalence}{equivalencia} y
colorea de rosa el matraz después de ella.

\textbf{6.} En la bureta; se lee en la parte inferior del menisco, con
el ojo a su altura, antes y después, a \qty{0.05}{mL}.

\textbf{7.} Amarillo pálido (\omterm{def:g9:ions:ion}{iones} hierro); cada gota violeta se
consume enseguida con los \omterm{def:g9:ions:ion}{iones} hierro(II) que quedan.

\textbf{8.} $0.0200 \times 0.0143 = \qty{2.86e-4}{mol}$.

\textbf{9.} $n(\ce{Fe^{2+}}) = 5 \times \num{2.86e-4} =
\qty{1.43e-3}{mol}$.

\textbf{10.} $\num{1.43e-3} \times 55.8 = \qty{0.0798}{g} =
\qty{79.8}{mg}$.

\textbf{11.} $(80 - 79.8)/80 \approx \qty{0.3}{\percent}$.

\textbf{12.} $\qty{0.1}{mL}$ de \qty{14.3}{mL} es un \qty{0.7}{\percent},
unos \qty{0.6}{mg}: el resultado, $79.8 \pm 0.6$ mg, concuerda con la
etiqueta.

\textbf{13.} $14.1$: $5 \times 0.0200 \times 0.0141 \times 55.8 =
\qty{78.7}{mg}$; $14.4$: \qty{80.4}{mg}.

\textbf{14.} $(79.8 + 78.7 + 80.4)/3 = \qty{79.6}{mg}$.

\textbf{15.} Los propios comprimidos difieren un poco en su contenido de
hierro, en alrededor de un \qty{1}{\percent}.

\textbf{16.} No: $V_E$ sería de unos \qty{2.9}{mL}, y el error de la
bureta, del \qty{3.5}{\percent}, cinco veces peor.

\textbf{17.} La vitamina C también se oxidaría con el permanganato: la
reacción ya no sería única, $V_E$ sería demasiado grande y el hierro se
sobrestimaría.

\textbf{18.} $\num{1.43e-3} \times 151.9 = \qty{0.217}{g}$, unos
\qty{217}{mg}.

\textbf{19.} $\num{1.43e-3} \times \num{6.02e23} = \num{8.6e20}$ \omterm{def:g9:ions:ion}{iones}.

\textbf{20.} Unos \qty{80}{mg} de hierro por comprimido: la etiqueta es
honrada.
\end{solution}
